\documentclass[UTF8,fontset=windows]{ctexart}

\usepackage[a4paper,margin=2.2cm]{geometry}
\usepackage{booktabs}
\usepackage{enumitem}
\usepackage{xcolor}
\usepackage{hyperref}
\usepackage{fvextra}

\hypersetup{
  colorlinks=true,
  linkcolor=blue,
  urlcolor=blue
}

\setlist{nosep}
\DefineVerbatimEnvironment{code}{Verbatim}{
  fontsize=\footnotesize,
  breaklines=true,
  breakanywhere=true,
  frame=single,
  rulecolor=\color{black!25}
}

\title{FG-MEET MATLAB 录屏演示脚本}
\author{}
\date{2026 年 6 月 25 日}

\begin{document}

\maketitle

\section{录屏目标}

本视频只演示 MATLAB 中的程序执行顺序，不讲解代码实现。目标是证明程序应从项目根目录运行，通过主入口读取 \texttt{txt} 输入文件并调用有限元核心求解，而不是直接运行 \texttt{materials} 文件夹下的辅助函数。

\section{录屏前准备}

\begin{enumerate}
  \item 打开 MATLAB。
  \item 确认交付包已经解压。
  \item 确认交付包根目录中包含：
  \begin{itemize}
    \item \texttt{setup\_paths.m}
    \item \texttt{run\_phase1\_static\_elastic.m}
    \item \texttt{run\_meet\_static.m}
    \item \texttt{cases/}
    \item \texttt{matlab/}
    \item \texttt{materials/}
    \item \texttt{output/}
  \end{itemize}
\end{enumerate}

\section{镜头 1：进入交付包根目录}

\textbf{画面操作：} MATLAB 命令行窗口。

\textbf{输入命令：}

\begin{code}
cd('C:\Users\You Huiwen\Desktop\fg-meet-workbench_code_delivery_2026-06-17\fg-meet-workbench_code_delivery_2026-06-17')
\end{code}

\textbf{画面展示：} 左侧当前文件夹切换到交付包根目录。

\textbf{旁白或字幕：}

\begin{quote}
首先进入交付包根目录，后续所有命令都从该目录执行。
\end{quote}

\section{镜头 2：初始化 MATLAB 路径}

\textbf{输入命令：}

\begin{code}
setup_paths
\end{code}

\textbf{预期现象：}

MATLAB 命令行输出工作目录、有限元核心目录等路径信息。

\textbf{旁白或字幕：}

\begin{quote}
先执行 \texttt{setup\_paths}，加载材料函数、工具函数和实验室有限元核心程序路径。
\end{quote}

\section{镜头 3：运行快速验证算例}

\textbf{输入命令：}

\begin{code}
run('run_phase1_static_elastic.m')
\end{code}

\textbf{预期现象：}

MATLAB 开始读取输入文件、装配矩阵并执行静力求解。运行完成后会输出结果文件路径。

\textbf{旁白或字幕：}

\begin{quote}
这是最小验证入口。该脚本会读取已有 \texttt{txt} 输入文件，并执行一次力载荷静力计算。
\end{quote}

\section{镜头 4：展示快速验证输出文件}

\textbf{输入命令：}

\begin{code}
dir output
\end{code}

\textbf{预期现象：}

在命令行中看到 \texttt{phase1\_static\_elastic\_summary.csv} 或相关输出文件。

\textbf{继续输入：}

\begin{code}
open('output/phase1_static_elastic_summary.csv')
\end{code}

\textbf{旁白或字幕：}

\begin{quote}
运行完成后，结果写入 \texttt{output} 文件夹。这里打开汇总表确认程序已经正常输出。
\end{quote}

\section{镜头 5：指定单个 txt 输入文件}

\textbf{输入命令：}

\begin{code}
paths = setup_paths;

caseFile = fullfile(paths.cases, ...
    'Thermal_CFFF_U_Vf0.6-30x30-10layer.txt');
\end{code}

\textbf{预期现象：}

工作区出现 \texttt{paths} 和 \texttt{caseFile} 变量。

\textbf{旁白或字幕：}

\begin{quote}
也可以手动指定某一个 \texttt{txt} 输入文件。该文件包含节点、单元、材料和边界条件。
\end{quote}

\section{镜头 6：运行力载荷工况}

\textbf{输入命令：}

\begin{code}
result_elastic = run_meet_static(caseFile, 'elastic', ...
    'OutTag', 'U_Vf06_elastic_demo');
\end{code}

\textbf{预期现象：}

MATLAB 执行力载荷静力求解，工作区出现 \texttt{result\_elastic}。

\textbf{旁白或字幕：}

\begin{quote}
这里运行力载荷工况，入口是 \texttt{run\_meet\_static}，不是材料辅助函数。
\end{quote}

\section{镜头 7：运行电载荷工况}

\textbf{输入命令：}

\begin{code}
result_electro = run_meet_static(caseFile, 'electro', ...
    'OutTag', 'U_Vf06_electro_demo');
\end{code}

\textbf{预期现象：}

MATLAB 执行电载荷工况，工作区出现 \texttt{result\_electro}。

\textbf{旁白或字幕：}

\begin{quote}
同一个输入文件可以切换为电载荷工况。
\end{quote}

\section{镜头 8：运行磁载荷工况}

\textbf{输入命令：}

\begin{code}
result_magneto = run_meet_static(caseFile, 'magneto', ...
    'OutTag', 'U_Vf06_magneto_demo');
\end{code}

\textbf{预期现象：}

MATLAB 执行磁载荷工况，工作区出现 \texttt{result\_magneto}。

\textbf{旁白或字幕：}

\begin{quote}
同一个输入文件也可以切换为磁载荷工况。
\end{quote}

\section{镜头 9：展示结果变量}

\textbf{输入命令：}

\begin{code}
result_elastic
result_electro
result_magneto
\end{code}

\textbf{预期现象：}

MATLAB 命令行展示结果结构体字段，例如中心挠度、温差均值、温差跨度等。

\textbf{旁白或字幕：}

\begin{quote}
这里展示三个工况的结果变量，说明三个主工况均已通过统一入口执行。
\end{quote}

\section{镜头 10：展示正确调用链路}

\textbf{画面操作：} 在 MATLAB 编辑器中打开 \texttt{run\_meet\_static.m}。

\textbf{需要展示的位置：}

\begin{itemize}
  \item \texttt{InputFile = caseFile}
  \item \texttt{Main\_FOSDLIN851T5MEET\_V4(...)}
\end{itemize}

\textbf{旁白或字幕：}

\begin{quote}
主入口会把指定的 \texttt{caseFile} 作为 \texttt{InputFile}，并调用实验室原有有限元核心程序完成装配和求解。
\end{quote}

\section{镜头 11：说明不应直接运行材料辅助函数}

\textbf{画面操作：} 打开 \texttt{materials/build\_layer\_materials.m}，只展示函数第一行。

\textbf{需要展示的位置：}

\begin{code}
function layers = build_layer_materials(nLayer, h, vf0, fgMode, varargin)
\end{code}

\textbf{旁白或字幕：}

\begin{quote}
\texttt{build\_layer\_materials} 是材料层参数辅助函数，需要输入参数。它不是主程序，不能零参数直接运行。
\end{quote}

\section{建议录屏结尾}

\textbf{字幕：}

\begin{quote}
正确执行顺序是先进入交付包根目录，再执行路径初始化和主程序入口。程序通过 \texttt{txt} 输入文件进入实验室有限元核心代码完成计算。
\end{quote}

\begin{code}
setup_paths
run('run_phase1_static_elastic.m')
run_meet_static(caseFile, loadCase)
\end{code}

\section{完整命令清单}

录屏时可以按以下命令顺序复制执行：

\begin{code}
cd('C:\Users\You Huiwen\Desktop\fg-meet-workbench_code_delivery_2026-06-17\fg-meet-workbench_code_delivery_2026-06-17')

setup_paths

run('run_phase1_static_elastic.m')

dir output

open('output/phase1_static_elastic_summary.csv')

paths = setup_paths;
caseFile = fullfile(paths.cases, ...
    'Thermal_CFFF_U_Vf0.6-30x30-10layer.txt');

result_elastic = run_meet_static(caseFile, 'elastic', ...
    'OutTag', 'U_Vf06_elastic_demo');

result_electro = run_meet_static(caseFile, 'electro', ...
    'OutTag', 'U_Vf06_electro_demo');

result_magneto = run_meet_static(caseFile, 'magneto', ...
    'OutTag', 'U_Vf06_magneto_demo');

result_elastic
result_electro
result_magneto
\end{code}

\end{document}
